% questions.tex —— 2004 年·普通高等学校招生考试·江苏卷 · 题目内容片段（无答案）
% 由 tools/generate_projects.py 依来源总账自 2004/2004_数学真题_江苏_试题.md 生成，勿手改（改源 MD 后重跑）

\gkpart{一、选择题：本大题共 12 小题，每小题 5 分，共 60 分．在每小题给出的四个选项中，只有一项是符合题目要求的．}

\begin{question}
\gkti{1}设集合 $P = \{1, 2, 3, 4\}$，$Q = \{x \mid |x| \leqslant 2, x \in \mathbf{R} \}$，则 $P \cap Q$ 等于\mbox{（　　）}

\begin{gkopts}(2)
\task A．$\{1, 2\}$
\task B．$\{3, 4\}$
\task C．$\{1\}$
\task D．$\{-2, -1, 0, 1, 2\}$
\end{gkopts}

\end{question}
\begin{question}
\gkti{2}函数 $y = 2\cos^2 x + 1$（$x \in \mathbf{R}$）的最小正周期为\mbox{（　　）}

\begin{gkopts}(4)
\task A．$\dfrac{\pi}{2}$
\task B．$\pi$
\task C．$2\pi$
\task D．$4\pi$
\end{gkopts}

\end{question}
\begin{question}
\gkti{3}从 4 名男生和 3 名女生中选出 4 人参加某个座谈会，若这 4 人中必须既有男生又有女生，则不同的选法共有\mbox{（　　）}

\begin{gkopts}(4)
\task A．$140$ 种
\task B．$120$ 种
\task C．$35$ 种
\task D．$34$ 种
\end{gkopts}

\end{question}
\begin{question}
\gkti{4}一平面截一球得到直径是 $6\text{ cm}$ 的圆面，球心到这个平面的距离是 $4\text{ cm}$，则该球的体积是\mbox{（　　）}

\begin{gkopts}(2)
\task A．$\dfrac{100\pi}{3}\text{ cm}^3$
\task B．$\dfrac{208\pi}{3}\text{ cm}^3$
\task C．$\dfrac{500\pi}{3}\text{ cm}^3$
\task D．$\dfrac{416\sqrt{3}\pi}{3}\text{ cm}^3$
\end{gkopts}

\end{question}
\begin{question}
\gkti{5}若双曲线 $\dfrac{x^2}{8} - \dfrac{y^2}{b^2} = 1$ 的一条准线与抛物线 $y^2 = 8x$ 的准线重合，则双曲线离心率为\mbox{（　　）}

\begin{gkopts}(4)
\task A．$\sqrt{2}$
\task B．$2\sqrt{2}$
\task C．$4$
\task D．$4\sqrt{2}$
\end{gkopts}

\Needspace{0.4\textheight}
\end{question}
\begin{question}
\gkti{6}某校为了了解学生的课外阅读情况，随机调查了 50 名学生，得到他们在某一天各自课外阅读所用时间的数据，结果用下面的条形图表示．根据条形图可得这 50 名学生这一天平均每人的课外阅读时间为\mbox{（　　）}
\nopagebreak[4]
\begin{center}\begin{tikzpicture}[>=stealth, xscale=1.5, yscale=0.18]
  % 网格虚线
  \draw[dashed, gray!50] (0, 5) -- (4.5, 5);
  \draw[dashed, gray!50] (0, 10) -- (4.5, 10);
  \draw[dashed, gray!50] (0, 15) -- (4.5, 15);
  \draw[dashed, gray!50] (0, 20) -- (4.5, 20);

  % 条形柱
  \draw[fill=gray!30] (-0.15, 0) rectangle (0.15, 5);
  \draw[fill=gray!30] (0.85, 0) rectangle (1.15, 20);
  \draw[fill=gray!30] (1.85, 0) rectangle (2.15, 10);
  \draw[fill=gray!30] (2.85, 0) rectangle (3.15, 10);
  \draw[fill=gray!30] (3.85, 0) rectangle (4.15, 5);

  % 坐标轴
  \draw[->, thick] (-0.3, 0) -- (4.8, 0) node[right] {\small 时间\,(小时)};
  \draw[->, thick] (0, 0) -- (0, 23) node[above] {\small 人数\,(人)};

  % 纵轴刻度
  \node[left] at (0, 0) {\small $0$};
  \node[left] at (0, 5) {\small $5$};
  \node[left] at (0, 10) {\small $10$};
  \node[left] at (0, 15) {\small $15$};
  \node[left] at (0, 20) {\small $20$};

  % 横轴刻度与标签
  \node[below] at (0, 0) {\small $0$};
  \node[below] at (1, 0) {\small $0.5$};
  \node[below] at (2, 0) {\small $1.0$};
  \node[below] at (3, 0) {\small $1.5$};
  \node[below] at (4, 0) {\small $2.0$};
\end{tikzpicture}
\end{center}
\nopagebreak[4]
\begin{gkopts}(4)
\task A．$0.6$ 小时
\task B．$0.9$ 小时
\task C．$1.0$ 小时
\task D．$1.5$ 小时
\end{gkopts}

\end{question}
\begin{question}
\gkti{7}$(2x + \sqrt{x})^4$ 的展开式中 $x^3$ 的系数是\mbox{（　　）}

\begin{gkopts}(4)
\task A．$6$
\task B．$12$
\task C．$24$
\task D．$48$
\end{gkopts}

\end{question}
\begin{question}
\gkti{8}若函数 $y = \log_a(x + b)$（$a > 0, a \neq 1$）的图象过两点 $(-1, 0)$ 和 $(0, 1)$，则\mbox{（　　）}

\begin{gkopts}(2)
\task A．$a = 2, b = 2$
\task B．$a = \sqrt{2}, b = 2$
\task C．$a = 2, b = 1$
\task D．$a = \sqrt{2}, b = \sqrt{2}$
\end{gkopts}

\end{question}
\begin{question}
\gkti{9}将一颗质地均匀的骰子（它是一种各面上分别标有点数 1，2，3，4，5，6 的正方体玩具）先后抛掷 3 次，至少出现一次 6 点的概率是\mbox{（　　）}

\begin{gkopts}(2)
\task A．$\dfrac{5}{216}$
\task B．$\dfrac{25}{216}$
\task C．$\dfrac{31}{216}$
\task D．$\dfrac{91}{216}$
\end{gkopts}

\end{question}
\begin{question}
\gkti{10}函数 $f(x) = x^3 - 3x + 1$ 在闭区间 $[-3, 0]$ 上的最大值、最小值分别是\mbox{（　　）}

\begin{gkopts}(4)
\task A．$1, -1$
\task B．$1, -17$
\task C．$3, -17$
\task D．$9, -19$
\end{gkopts}

\end{question}
\begin{question}
\gkti{11}设 $k > 1, f(x) = k(x - 1)$（$x \in \mathbf{R}$）．在平面直角坐标系 $xOy$ 中，函数 $y = f(x)$ 的图象与 $x$ 轴交于 $A$ 点，它的反函数 $y = f^{-1}(x)$ 的图象与 $y$ 轴交于 $B$ 点，并且这两个函数的图象交于 $P$ 点．已知四边形 $OAPB$ 的面积是 3，则 $k$ 等于\mbox{（　　）}

\begin{gkopts}(4)
\task A．$3$
\task B．$\dfrac{3}{2}$
\task C．$\dfrac{4}{3}$
\task D．$\dfrac{6}{5}$
\end{gkopts}

\end{question}
\begin{question}
\gkti{12}设函数 $f(x) = -\dfrac{x}{1 + |x|}$（$x \in \mathbf{R}$），区间 $M = [a, b]$（$a < b$），集合 $N = \{y \mid y = f(x), x \in M\}$，则使 $M = N$ 成立的实数对 $(a, b)$ 有\mbox{（　　）}

\begin{gkopts}(4)
\task A．$0$ 个
\task B．$1$ 个
\task C．$2$ 个
\task D．无数多个
\end{gkopts}

\end{question}
\gkpart{二、填空题：本大题共 4 小题，每小题 4 分，共 16 分．把答案填在题中横线上．}

\begin{question}
\gkti{13}二次函数 $y = ax^2 + bx + c$（$x \in \mathbf{R}$）的部分对应值如下表：

\begin{center}
\begingroup\setlength{\tabcolsep}{4pt}
\begin{tabular}{ccccccccc}
\toprule
$x$ & $-3$ & $-2$ & $-1$ & $0$ & $1$ & $2$ & $3$ & $4$ \\
\midrule
$y$ & $6$ & $0$ & $-4$ & $-6$ & $-6$ & $-4$ & $0$ & $6$ \\
\bottomrule
\end{tabular}
\endgroup
\end{center}

则不等式 $ax^2 + bx + c > 0$ 的解集是 \underline{\hspace{3cm}}．

\end{question}
\begin{question}
\gkti{14}以点 $(1, 2)$ 为圆心，与直线 $4x + 3y - 35 = 0$ 相切的圆的方程是 \underline{\hspace{3cm}}．

\end{question}
\begin{question}
\gkti{15}设数列 $\{a_n\}$ 的前 $n$ 项和为 $S_n$，$S_n = \dfrac{a_1(3^n - 1)}{2}$（对于所有 $n \geqslant 1$），且 $a_4 = 54$，则 $a_1$ 的数值是 \underline{\hspace{3cm}}．

\end{question}
\begin{question}
\gkti{16}平面向量 $\vec{a}, \vec{b}$ 中，已知 $\vec{a} = (4, -3)$，$|\vec{b}| = 1$，且 $\vec{a} \cdot \vec{b} = 5$，则向量 $\vec{b} =$ \underline{\hspace{3cm}}．

\end{question}
\gkpart{三、解答题：本大题共 6 小题，共 74 分．解答应写出文字说明、证明过程或演算步骤．}

\begin{question}
\gkti{17}\textbf{（本小题满分 12 分）} 

已知 $0 < \alpha < \dfrac{\pi}{2}$，$\tan\dfrac{\alpha}{2} + \cot\dfrac{\alpha}{2} = \dfrac{5}{2}$，求 $\sin\left(\alpha - \dfrac{\pi}{3}\right)$ 的值．

\Needspace{0.4\textheight}
\end{question}
\begin{question}
\gkti{18}\textbf{（本小题满分 12 分）} 

在棱长为 4 的正方体 $ABCD - A_1B_1C_1D_1$ 中，$O$ 是正方形 $A_1B_1C_1D_1$ 的中心，点 $P$ 在棱 $CC_1$ 上，且 $CC_1 = 4CP$．

（1）求直线 $AP$ 与平面 $BCC_1B_1$ 所成的角的大小（结果用反三角函数值表示）；

（2）设 $O$ 点在平面 $D_1AP$ 上的射影是 $H$，求证：$D_1H \perp AP$；

（3）求点 $P$ 到平面 $ABD_1$ 的距离．
\nopagebreak[4]
\begin{center}\begin{tikzpicture}[>=stealth, scale=1.0]
  % 顶点坐标
  \coordinate (D) at (0, 0);
  \coordinate (C) at (3.2, 0);
  \coordinate (A) at (-1.5, -1.2);
  \coordinate (B) at (1.7, -1.2);

  \coordinate (D1) at (0, 3.0);
  \coordinate (C1) at (3.2, 3.0);
  \coordinate (A1) at (-1.5, 1.8);
  \coordinate (B1) at (1.7, 1.8);

  % 点 P 在 CC1 上，CP = 1/4 CC1
  \coordinate (P) at ($(C)!0.25!(C1)$);

  % 上底面中心 O
  \coordinate (O) at ($(A1)!0.5!(C1)$);

  % 射影点 H (在平面 D1AP 上)
  \coordinate (H) at (0.9, 1.4);

  % 看不见的棱 (虚线)
  \draw[dashed, thick] (D) -- (A);
  \draw[dashed, thick] (D) -- (C);
  \draw[dashed, thick] (D) -- (D1);

  % 外轮廓与可见棱 (实线)
  \draw[thick] (A) -- (B) -- (C);
  \draw[thick] (A1) -- (B1) -- (C1) -- (D1) -- cycle;
  \draw[thick] (A) -- (A1);
  \draw[thick] (B) -- (B1);
  \draw[thick] (C) -- (C1);

  % 连线 AP, D1A, D1P
  \draw[thick, blue!80!black] (D1) -- (A);
  \draw[thick, blue!80!black] (A) -- (P);
  \draw[thick, blue!80!black] (D1) -- (P);

  % 辅助线 D1O, OH, D1H
  \draw[thick] (D1) -- (O);
  \draw[dashed] (O) -- (H);
  \draw[thick, red!70!black] (D1) -- (H);

  % 标记点
  \fill (O) circle (1.5pt) node[above right] {$O$};
  \fill (P) circle (1.5pt) node[right] {$P$};
  \fill (H) circle (1.5pt) node[below right] {$H$};

  % 顶点标签
  \node[left] at (A) {$A$};
  \node[below right] at (B) {$B$};
  \node[right] at (C) {$C$};
  \node[above left] at (D) {$D$};
  \node[left] at (A1) {$A_1$};
  \node[below right] at (B1) {$B_1$};
  \node[above right] at (C1) {$C_1$};
  \node[above left] at (D1) {$D_1$};
\end{tikzpicture}
\end{center}
\nopagebreak[4]
\end{question}
\begin{question}
\gkti{19}\textbf{（本小题满分 12 分）} 

制定投资计划时，不仅要考虑可能获得的盈利，而且要考虑可能出现的亏损．某投资人打算投资甲、乙两个项目．根据预测，甲、乙项目可能的最大盈利率分别为 $100\%$ 和 $50\%$，可能的最大亏损分别为 $30\%$ 和 $10\%$．投资人计划投资金额不超过 10 万元，要求确保可能的资金亏损不超过 1.8 万元．问投资人对甲、乙两个项目各投资多少万元，才能使可能的盈利最大？

\end{question}
\begin{question}
\gkti{20}\textbf{（本小题满分 12 分）} 

设无穷等差数列 $\{a_n\}$ 的前 $n$ 项和为 $S_n$．

（1）若首项 $a_1 = \dfrac{3}{2}$，公差 $d = 1$，求满足 $S_{k^2} = (S_k)^2$ 的正整数 $k$；

（2）求所有的无穷等差数列 $\{a_n\}$，使得对于一切正整数 $k$ 都有 $S_{k^2} = (S_k)^2$ 成立．

\end{question}
\begin{question}
\gkti{21}\textbf{（本小题满分 12 分）} 

已知椭圆的中心在原点，离心率为 $\dfrac{1}{2}$，一个焦点是 $F(-m, 0)$（$m$ 是大于 0 的常数）．

（1）求椭圆的方程；

（2）设 $Q$ 是椭圆上的一点，且过点 $F, Q$ 的直线 $l$ 与 $y$ 轴交于点 $M$．若 $\vec{MQ} = 2\vec{QF}$，求直线 $l$ 的斜率．

\end{question}
\begin{question}
\gkti{22}\textbf{（本小题满分 14 分）} 

\Needspace{5\baselineskip}
已知函数 $f(x)$（$x \in \mathbf{R}$）满足下列条件：对任意的实数 $x_1, x_2$ 都有
\nopagebreak[4]
\[
\lambda(x_1 - x_2)^2 \leqslant (x_1 - x_2)[f(x_1) - f(x_2)] \quad \text{\text{和}} \quad |f(x_1) - f(x_2)| \leqslant |x_1 - x_2|,
\]
其中 $\lambda$ 是大于 0 的常数．设实数 $a_0, a, b$ 满足 $f(a_0) = 0$ 和 $b = a - \lambda f(a)$．

（1）证明：$\lambda \leqslant 1$，并且不存在 $b_0 \neq a_0$，使得 $f(b_0) = 0$；

（2）证明：$(b - a_0)^2 \leqslant (1 - \lambda^2)(a - a_0)^2$；

（3）证明：$[f(b)]^2 \leqslant (1 - \lambda^2)[f(a)]^2$．

\end{question}